% !PRINTED 297
\weberpage{297}
%<<CONT
\weq{4}{
\begin{aligned}
\frac{d\varphi}{d\vartheta} &= -n R^n \sin n\vartheta - (n-1) p_1 R^{n-1} \sin[(n-1)\vartheta + q_1] \\
&\quad - \cdots - p_{n-1} \sin(\vartheta + q_{n-1}) \\[4pt]
\frac{d\psi}{d\vartheta} &= n R^n \cos n\vartheta + (n-1) p_1 R^{n-1} \cos[(n-1)\vartheta + q_1)] \\
&\quad + \cdots + p_{n-1} \cos(\vartheta + q_{n-1}).
\end{aligned}
} % [sic] stray closing parenthesis after q_1 in the cos-bracket of the second line, as printed

We now divide the periphery of the circle of radius $R$ into
$4n$ parts of the angular magnitude
\[
2\omega = 2\,\frac{\pi}{4n},
\]
beginning at $\vartheta = -\omega$. The points of division are therefore
\[
-\omega,\ \omega,\ 3\omega,\ 5\omega,\ \dots (8n - 3)\omega,
\]
and let the intervals be denoted in order by $1, 2, 3 \dots 4n$.
Fig. 15 shows this division for $n = 3$.

\weberfigureimage{weber-b1-fig15.png}{0.54}

In the intervals
\[
1, 3, 5 \dots (4n - 1)
\]
$\cos n\vartheta$, taken absolutely, is greater than $1 : \sqrt{2}$, and
is alternately of positive and negative sign. In the
intervals
\[
2, 4, 6 \dots 4n
\]
$\sin n\vartheta$ is absolutely greater than $1 : \sqrt{2}$ and likewise
alternately of positive and negative sign.

But $R$ may be taken so large that in the intervals
$1, 3, 5 \dots (4n - 1)$ the sign of $\varphi$ and of $\dfrac{d\psi}{d\vartheta}$ agrees with that
of $\cos n\vartheta$, and in the intervals $2, 4, 6 \dots 4n$ the sign
of $\psi$ and of $\dfrac{d\varphi}{d\vartheta}$ agrees with that of $\sin n\vartheta$.

Since, then, $\varphi$ changes its sign in the intervals $2, 4, 6 \dots 4n$,
it must pass through zero in each of these intervals;
and since the derivative $\dfrac{d\varphi}{d\vartheta}$ does not change its sign, it can
pass through zero only \gesperrt{once} in each.
